\clearpage\SourceLeaf{186}
\begin{center}\textsc{Gauss} to \textsc{Gerling}. Göttingen, 25 August 1818.\end{center}

.\ .\ .\ .\ .\ .\ Of your questions about \textsc{Lorenz}'s textbook I can
answer only a few today. .\ .\ .\ .\ .\ .

.\ .\ .\ .\ .\ .\ I am glad you have the courage to express yourself as
if you acknowledged the possibility that our theory of parallels, and hence our entire
geometry, might be false. But the wasps whose nest you
are disturbing will fly about your head. .\ .\ .\ .\ .\ .

\begin{center}\rule{3cm}{0.4pt}\end{center}

\begin{center}[\textsc{Gerling} \emph{to} \textsc{Gauss}. \emph{Marburg, 25 January 1819.}]\end{center}

\{.\ .\ .\ .\ .\ .\ .\ .\ I have now worded the passage on the theory of parallels as follows.
»The proposition in \S\ 72 is set out in Euclid's Elements (Book 1, Axiom 11) as an
axiom. That it is no axiom but requires a
proof can be seen from the fact that it contains two essentially different
conceptions (the angles of two lines with a third and their meeting
one another), whose \textls{necessary} connection must be established by
a proof. --- This proof (the
theory of parallels) has been attempted in various ways by acute
mathematicians, but none has yet been found that is entirely
satisfactory. As long as it is lacking, the proposition and all that
rests on it remain a hypothesis whose validity for our lives is indeed
shown by experience, but whose \textls{general, necessary}
truth could be doubted without absurdity«.

Ad vocem the theory of parallels, I must tell you something else and discharge a
commission. Last year I learned that my colleague
\textsc{Schweikart} (Prof. juris, now Prorector) used to occupy himself a great deal with
mathematics and had written about parallels in particular.
I therefore asked him to lend me his book. In promising to do so,
he told me that he now saw that errors had occurred in his book
(1808) (he had, for\,example, used quadrilaterals with equal angles as
a primitive concept), but that he had not ceased
to occupy himself with the subject,
